%%%
%%% Radical "⿇"
%%%
\section*{Radical 200: ``⿇''}\addcontentsline{toc}{section}{Radical 200: ⿇}\addcontentsline{loh}{figure}{\#\#\#\# 200: ⿇}

%%%%%%%%%% 麻 %%%%%%%%%%
\subsection*{麻}\addcontentsline{loh}{figure}{麻}

\begin{Entry}{麻}{11}{⿇}[Kangxi 200]
  \begin{Phonetics}{麻}{ma1}
    \definition{adj.}{sombrio; escuro; completamente escuro}
  \end{Phonetics}
  \begin{Phonetics}{麻}{ma2}[][HSK 7-9]
    \definition*{s.}{Sobrenome: Ma}
    \definition{adj.}{áspero; grosseiro | marcado; manchado | espinhas; manchas ásperas; cicatrizes deixadas após a varíola}
    \definition[棵,株]{s.}{nome geral para cânhamo, linho, etc. | fibra de cânhamo, linho, etc. para têxteis | sésamo; gergelim | marcas de varíola; um rosto com marcas de varíola}
    \definition{v.}{anestesiar | corromper (a mente de alguém); envenenar}
  \end{Phonetics}
\end{Entry}

\begin{Entry}{麻木}{11,4}{⿇,⽊}
  \begin{Phonetics}{麻木}{ma2mu4}[][HSK 7-9]
    \definition{adj.}{dormente; descreve uma sensação de dormência ou perda de sensibilidade em uma parte do corpo devido ao frio ou à inatividade prolongada | entorpecido; insensível; sem vida; apático; de raciocínio lento}
  \end{Phonetics}
\end{Entry}

\begin{Entry}{麻将}{11,9}{⿇,⼨}
  \begin{Phonetics}{麻将}{ma2jiang4}[][HSK 7-9]
    \definition*[副]{s.}{Mahjong}
  \end{Phonetics}
\end{Entry}

\begin{Entry}{麻烦}{11,10}{⿇,⽕}
  \begin{Phonetics}{麻烦}{ma2fan5}[][HSK 3]
    \definition{adj.}{incômodo; inconveniente; complicado; trabalhoso; burocrático | incômodo; inconveniente; (a situação) é confusa e complicada}
    \definition[个,些,点,堆]{s.}{problema; inconveniência; assuntos complicados e difíceis de resolver}
    \definition{v.}{incomodar; perturbar; incomodar alguém; irritar; aborrecer; causar incômodo ou sobrecarregar outras pessoas}
  \synonymref{烦恼}{fan2nao3}
  \synonymref{艰难}{jian1nan2}
  \synonymref{困难}{kun4nan5}
  \synonymref{难处}{nan2chu4}
  \synonymref{难处}{nan2chu5}
  \antonymref{便利}{bian4li4}
  \antonymref{方便}{fang1bian4}
  \antonymref{简便}{jian3bian4}
  \antonymref{简易}{jian3yi4}
  \antonymref{容易}{rong2yi4}
  \antonymref{顺利}{shun4li4}
  \end{Phonetics}
\end{Entry}

\begin{Entry}{麻痹}{11,13}{⿇,⽧}
  \begin{Phonetics}{麻痹}{ma2bi4}[][HSK 7-9]
    \definition{adj.}{insensível; descuidado; negligente; descreve um estado de entorpecimento mental e perda de vigilância}
    \definition{v.}{ficar paralisado; ficar insensível; perder total ou parcialmente as funções sensoriais e motoras em uma parte do corpo | estar insensível; baixar a guarda; entorpecer a mente}
  \end{Phonetics}
\end{Entry}

\begin{Entry}{麻辣}{11,14}{⿇,⾟}
  \begin{Phonetics}{麻辣}{ma2la4}[][HSK 7-9]
    \definition{adj.}{picante; quente e dormente; descreve um sabor como dormente e picante, semelhante à pimenta"-de"-Sichuan ou à pimenta malagueta}
  \end{Phonetics}
\end{Entry}

\begin{Entry}{麻辣豆腐}{11,14,7,14}{⿇,⾟,⾖,⾁}
  \begin{Phonetics}{麻辣豆腐}{ma2la4 dou4fu5}
    \definition{s.}{tofú guisado em molho picante (prato)}
  \end{Phonetics}
\end{Entry}

\begin{Entry}{麻醉}{11,15}{⿇,⾣}
  \begin{Phonetics}{麻醉}{ma2zui4}[][HSK 7-9]
    \definition{v.}{narcotizar; anestesiar; utilizar drogas, acupuntura ou outros métodos para deixar temporariamente todo ou parte de um organismo inconsciente | corromper (a mente de alguém); desgastar (a força de vontade de alguém); descreve o uso de certos métodos para desgastar a vontade de uma pessoa, fazendo com que ela perca a capacidade de distinguir o certo do errado}
  \end{Phonetics}
\end{Entry}

%%%%% EOF %%%%%

